%-----------------------------------
% Anhang II: Kurzform der verwendeten Prompts
% Leitfaden Abschnitt 6.2.4: "Daher soll eine Kurzform der Prompts (ohne
% Antwort) in den Anhang aufgenommen werden."
%
% ZU PRUEFEN: Die Eintraege beschreiben die Art der Nutzung. Datum und Wortlaut
% sind gegen die tatsaechlich gestellten Eingaben abzugleichen. Prompts, die
% hier nicht abgebildet sind, sind zu ergaenzen.
%-----------------------------------
\section{Kurzform der verwendeten Prompts}
\label{anhangPrompts}

\autoref{tab:prompts} gibt die an ein Sprachmodell gerichteten Eingaben in gek\"urzter Form wieder. Die Antworten des Modells sind nicht aufgef\"uhrt. Die Eingaben zielten durchg\"angig auf Vorschl\"age und Hinweise; die Auswahl der Quellen, die inhaltlichen Entscheidungen und die endg\"ultigen Formulierungen erfolgten durch den Verfasser.

\begin{table}[htbp]
	\centering
	\caption{Kurzform der verwendeten Prompts}
	\label{tab:prompts}
	\begin{tabularx}{\textwidth}{@{}p{2.15cm}p{2.4cm}X@{}}
		\toprule
		Datum & Bezug & Prompt (Kurzform) \\
		\midrule
		\multicolumn{3}{@{}l}{\textit{Unterst\"utzung bei der Quellensuche}} \\
		\addlinespace[2pt]
		$\bullet$ & Kap.~2.2.3 & Welche begutachteten Ver\"offentlichungen behandeln die Klassifikation von Service-Desk-Tickets mit Sprachmodellen? Nur Titel und Fundstelle nennen, keine Inhaltsangabe. \\
		$\bullet$ & Kap.~4.2.3 & Suchbegriffe f\"ur die Recherche zu Retrieval-Augmented Generation im IT-Support vorschlagen, die ich in IEEE Xplore und der ACM Digital Library einsetzen kann. \\
		$\bullet$ & Kap.~5.1.2 & Diese Liste gefundener Beitr\"age danach ordnen, ob sie sich auf Endanwender-Support oder auf den Betrieb von Cloud-Diensten beziehen. \\
		\addlinespace[4pt]
		\multicolumn{3}{@{}l}{\textit{Unterst\"utzung beim Formulieren}} \\
		\addlinespace[2pt]
		$\bullet$ & Kap.~2.1 & F\"ur den \"Ubergang von 2.1.1 zu 2.1.2 zwei bis drei Varianten eines \"Uberleitungssatzes vorschlagen, aus denen ich ausw\"ahle. \\
		$\bullet$ & Kap.~3.3 & Diesen von mir geschriebenen Absatz auf Verst\"andlichkeit und Satzl\"ange pr\"ufen und Stellen benennen, an denen er unklar wird. \\
		$\bullet$ & Kap.~4.4 & Pr\"ufen, ob meine Argumentation zum Fallback in sich schl\"ussig ist, und offene Stellen benennen, ohne den Text neu zu schreiben. \\
		\bottomrule
	\end{tabularx}
\end{table}
